%============================================================
% CH2 · PATTERN 4 — EXERCISES, HINTS, ANSWERS
%============================================================
\subsection{Pattern 4: Exercises and Going Further}
\label{ssec:ch2-p4-ex}

\begin{patternbox}{Pattern 4 Summary Card}
\renewcommand{\arraystretch}{1.25}
\begin{tabularx}{\linewidth}{@{}>{\bfseries}l X@{}}
Recognise & ``all paths'', ``the best path'', ``each vertex / cell at most once'' on a \emph{tiny} input ($n \le 10$--$15$, $\le 25$ cells, a word of $\le 15$ letters), or a capped count \\
Core & DFS whose marks are removed on exit: \texttt{onPath[]} $=$ the current path; one call per simple path \\
State & everything the answer needs, kept on the path: add on enter, undo on exit (path list, sum, colour counts, prefix sums) \\
Cost & (number of simple paths) $\times$ degree; count it on the worst input before coding \\
Cut & bounds (value so far $+$ reachable rest), dead ends (unreachable target, forced cells, split regions), parity, caps \\
Linear cases & DAG: no marks, memoise on the last vertex (DP); tree: one path per vertex, the path stack is a plain DFS \\
Top traps & forgetting to un-mark or to restore; going on past the target; plain DFS marks in a backtracking; a prune without proof; deep recursion on $10^5$ \\
\end{tabularx}
\end{patternbox}

\exlegend
\begin{enumerate}[leftmargin=*, label=\textbf{E4.\arabic*}, itemsep=6pt]
    \item \exeasy\ \pclassic \exprob{One-stroke Path}{AtCoder ABC054 C}{https://atcoder.jp/contests/abc054/tasks/abc054_c}. A simple undirected graph with $N \le 8$ vertices. Count the paths that start at vertex 1 and visit every vertex exactly once.
    \item \exmedium\ \pclassic \lc{Medium}\ \exprob{Word Search}{LeetCode 79}{https://leetcode.com/problems/word-search/}. An $m \times n$ board of letters ($m, n \le 6$) and a word of length $\le 15$. Can the word be read along side-adjacent cells, each cell used at most once?
    \item \exeasy\ \pconcept \exprob{Remembering the Days}{AtCoder ABC317 C}{https://atcoder.jp/contests/abc317/tasks/abc317_c}. $N \le 10$ towns, roads with lengths up to $10^8$. Start anywhere and travel without visiting a town twice; maximise the total length of the roads used.
    \item \exmedium\ \pconcept \exprob{Count Simple Paths}{AtCoder ABC284 E}{https://atcoder.jp/contests/abc284/tasks/abc284_e}. A simple undirected graph with $N, M \le 2\cdot10^5$ and every degree $\le 10$. Let $K$ be the number of simple paths starting at vertex 1 (the path ``1'' included). Print $\min(K, 10^6)$.
    \item \exmedium\ \pconcept \exprob{Unique Color}{AtCoder ABC198 E}{https://atcoder.jp/contests/abc198/tasks/abc198_e}. A tree with $N \le 10^5$ vertices; vertex $i$ has colour $C_i$. Vertex $x$ is good if no other vertex on the path from 1 to $x$ has the colour of $x$. Print all good vertices.
    \item \exhard\ \pconcept \cf{1700}\ \exprob{Path Prefixes}{Codeforces 1714G}{https://codeforces.com/problemset/problem/1714/G}. A tree rooted at 1 ($\sum n \le 2\cdot10^5$); every edge has two costs $a_j, b_j \ge 1$. For every vertex $i \ge 2$ let $A_i$ be the sum of $a$ on the root path; $r_i$ is the length of the longest prefix of that path whose sum of $b$ is at most $A_i$. Print $r_2, \dots, r_n$.
    \item \exhard\ \ptricky \exprob{Grid Paths}{CSES 1625}{https://cses.fi/problemset/task/1625}. Count the paths of 48 steps from the upper-left to the lower-left square of a $7 \times 7$ grid that visit all 49 squares, matching a 48-character description over \texttt{D}, \texttt{U}, \texttt{L}, \texttt{R} and \texttt{?} (any direction). There are $88\,418$ such paths in total; the time limit is 1 second.
\end{enumerate}

\subsubsection*{Hints}
\begin{enumerate}[leftmargin=*, label=\textbf{E4.\arabic*}]
    \item $7! = 5040$ orders at most. Template 1 without a target: count a path when its length reaches $N$.
    \item Backtrack letter by letter; mark a cell by overwriting it. Before searching at all, compare letter counts.
    \item Template 2 with edge weights; every node of the search tree is a candidate, and every town a start.
    \item Every call of the backtracking is one new simple path. What does the cap do to the number of calls? How deep can the recursion become?
    \item The path from 1 to $x$ in a tree is the DFS stack when $x$ is entered. Keep a count per colour.
    \item On the stack, keep the prefix sums of $b$ along the root path. They are increasing.
    \item Two prunes are enough: (1) moving straight into a blocked square while both side squares are free splits the free squares; (2) a free neighbour with at most one other free neighbour must be entered right now.
\end{enumerate}

\subsubsection*{Answers}

\begin{solutionbox}[E4.1 One-stroke Path]
Backtracking from vertex 1 with \texttt{onPath} marks; a path is complete when it has $N$ vertices. At most $(N-1)! = 5040$ complete paths, and at most $13\,700$ calls in total (the number of simple paths from one vertex of $K_8$, from the table of the pattern theory).
\begin{center}
\begin{tikzpicture}[x=1.1cm, y=0.9cm]
    \tikzset{gv/.append style={minimum size=6.5mm, font=\scriptsize\bfseries}}
    \node[gvsrc] (v1) at (0,0) {1}; \node[gv] (v3) at (1.2,0) {3}; \node[gv] (v4) at (2.4,0) {4};
    \node[gv] (v5) at (3.4,0.8) {5}; \node[gv] (v6) at (3.4,-0.8) {6}; \node[gv] (v7) at (4.8,-0.8) {7}; \node[gvdst] (v2) at (4.8,0.6) {2};
    \foreach \a/\b in {4/6} \draw[ge] (v\a) -- (v\b);
    \foreach \a/\b in {1/3, 3/4, 4/5, 5/6, 6/7, 7/2} \draw[gpath] (v\a) -- (v\b);
    \node[note, text width=4.4cm, anchor=west] at (5.4,0) {Sample 2: the only one-stroke path from 1 is $1, 3, 4, 5, 6, 7, 2$ (the edge $4{-}6$ is unused).};
\end{tikzpicture}
\end{center}
\begin{lstlisting}[style=cpp]
#include <bits/stdc++.h>
using namespace std;

// One-stroke Path (ABC054 C): count the paths that start at 1 and visit every vertex exactly once.
int n, m, answer = 0;
bool adj[9][9];
bool onPath[9];

void go(int u, int visited) {                // visited = number of vertices on the path
    if (visited == n) { answer++; return; }  // the path covers everything: one more Hamiltonian path
    for (int v = 1; v <= n; v++)
        if (adj[u][v] && !onPath[v]) {
            onPath[v] = true;                // enter v
            go(v, visited + 1);
            onPath[v] = false;               // leave v: other paths may use it
        }
}

int main() {
    cin >> n >> m;
    for (int i = 0; i < m; i++) {
        int a, b;
        cin >> a >> b;
        adj[a][b] = adj[b][a] = true;
    }
    onPath[1] = true;
    go(1, 1);
    cout << answer << "\n";
}
\end{lstlisting}
Samples: \textbf{2} (the triangle: $1,2,3$ and $1,3,2$) and \textbf{1}. Checked against a brute force over all $(N-1)!$ orders of $2..N$ (\texttt{next\_permutation}) on 500 random graphs with $N \le 8$: identical.
\end{solutionbox}

\begin{solutionbox}[E4.2 Word Search]
Backtracking over cells: \texttt{match(i, r, c)} checks that cell $(r,c)$ holds \texttt{word[i]}, overwrites it with \texttt{'\#'}, tries the four neighbours for \texttt{word[i+1]} and restores the letter. Two prunes, both provable: if the board has fewer copies of some letter than the word needs, the answer is \texttt{false} without any search; and a path read backwards is still a path, so the word may be searched from its rarer end.
\begin{lstlisting}[style=cpp]
#include <bits/stdc++.h>
using namespace std;

// Word Search (LeetCode 79): backtracking over cells; a used cell is marked '#' and restored on exit.
class Solution {
    int R, C;
    bool match(vector<vector<char>>& b, const string& w, int i, int r, int c) {
        if (r < 0 || r >= R || c < 0 || c >= C || b[r][c] != w[i]) return false;
        if (i + 1 == (int)w.size()) return true;
        char keep = b[r][c];
        b[r][c] = '#';                                       // enter: the cell is on the path
        bool ok = match(b, w, i + 1, r + 1, c) || match(b, w, i + 1, r - 1, c) ||
                  match(b, w, i + 1, r, c + 1) || match(b, w, i + 1, r, c - 1);
        b[r][c] = keep;                                      // exit: restore it
        return ok;
    }
public:
    bool exist(vector<vector<char>>& board, string word) {
        R = board.size(); C = board[0].size();
        map<char, int> cnt;
        for (auto& row : board) for (char ch : row) cnt[ch]++;
        map<char, int> need = cnt;
        for (char ch : word) if (--need[ch] < 0) return false;   // prune 1: not enough letters
        if (cnt[word.back()] < cnt[word[0]])                     // prune 2: start at the rarer end
            reverse(word.begin(), word.end());
        for (int r = 0; r < R; r++)
            for (int c = 0; c < C; c++)
                if (match(board, word, 0, r, c)) return true;
        return false;
    }
};

int main() {
    vector<vector<char>> b = {{'A', 'B', 'C', 'E'}, {'S', 'F', 'C', 'S'}, {'A', 'D', 'E', 'E'}};
    Solution s;
    for (string w : {"ABCCED", "SEE", "ABCB"}) cout << boolalpha << s.exist(b, w) << " ";
    cout << "\n";                                                // true true false
}
\end{lstlisting}
Examples: \textbf{true, true, false}. Measured on a $6 \times 6$ board of \texttt{A}s and the word \texttt{AAAAAAAAAAAAAAB}: without prune 1 the search makes $9\,030\,868$ calls before answering \texttt{false}; with it, none. Checked against a brute force over all ordered tuples of distinct cells (boards up to $3 \times 3$, words up to 5 letters over \texttt{AB} or \texttt{ABC}) on 3000 random cases, 1056 of them \texttt{true}: identical, and the board was restored every time.
\end{solutionbox}

\begin{solutionbox}[E4.3 Remembering the Days]
Template 2 with edge weights: \texttt{go(u, len)} updates the answer at every node and extends the path along every road to an unused town. All $N$ towns are tried as starts. $N \le 10$ is small enough even for the complete graph: $K_{10}$ has $986\,410$ simple paths from each town (the table of the pattern theory), so $10$ starts make about $10^7$ calls --- under half a second here. Lengths up to $10^8$ times 9 roads need \texttt{long long}.
\begin{lstlisting}[style=cpp]
#include <bits/stdc++.h>
using namespace std;
typedef long long ll;

// Remembering the Days (ABC317 C): maximum total length of a simple path, any start, N <= 10.
int n, m;
ll w[11][11];                    // 0 = no road
bool onPath[11];
ll best = 0;

void go(int u, ll len) {
    best = max(best, len);       // every simple path is a candidate, wherever it stops
    for (int v = 1; v <= n; v++)
        if (w[u][v] && !onPath[v]) {
            onPath[v] = true;
            go(v, len + w[u][v]);
            onPath[v] = false;
        }
}

int main() {
    cin >> n >> m;
    for (int i = 0; i < m; i++) {
        int a, b;
        ll c;
        cin >> a >> b >> c;
        w[a][b] = w[b][a] = c;
    }
    for (int s = 1; s <= n; s++) {
        onPath[s] = true;
        go(s, 0);
        onPath[s] = false;
    }
    cout << best << "\n";
}
\end{lstlisting}
Samples: \textbf{1110} ($4 \to 1 \to 3 \to 2$), \textbf{1}, \textbf{20}. Checked against a brute force over all $N!$ orders of the towns (best prefix that follows roads) on 500 random graphs with $N \le 8$: identical.
\end{solutionbox}

\begin{solutionbox}[E4.4 Count Simple Paths]
Every call of the backtracking enters a vertex along a new simple path, so ``count the calls and stop at $10^6$'' needs at most $10^6$ calls, each scanning $\le 10$ neighbours. The trap is that the obvious recursive code can still die: a path graph $1 - 2 - \dots - 2\cdot10^5$ has only $2\cdot10^5$ simple paths from 1 (below the cap) but makes the recursion $2\cdot10^5$ deep. The iterative version below keeps one iterator per vertex and \textbf{resets it when the vertex is left}, because the vertex may be entered again on another path.
\begin{lstlisting}[style=cpp]
#include <bits/stdc++.h>
using namespace std;

// Count Simple Paths (ABC284 E): number of simple paths from vertex 1, capped at 10^6.
// Every call of the backtracking IS one path, so the cap bounds the number of calls.
// Iterative: a path can be 2*10^5 vertices long.
int main() {
    ios::sync_with_stdio(false);
    cin.tie(nullptr);
    const long long CAP = 1000000;
    int n, m;
    cin >> n >> m;
    vector<vector<int>> adj(n + 1);
    for (int i = 0; i < m; i++) {
        int u, v;
        cin >> u >> v;
        adj[u].push_back(v);
        adj[v].push_back(u);
    }
    vector<char> onPath(n + 1, 0);
    vector<size_t> it(n + 1, 0);
    vector<int> st;
    long long count = 1;                       // the path "1" of length 0
    st.push_back(1); onPath[1] = 1;
    while (!st.empty() && count < CAP) {
        int u = st.back();
        if (it[u] == adj[u].size()) {          // u has no more extensions: leave it
            onPath[u] = 0; it[u] = 0;          // reset its iterator: u may be entered again later
            st.pop_back();
            continue;
        }
        int v = adj[u][it[u]++];
        if (onPath[v]) continue;
        onPath[v] = 1; st.push_back(v);        // enter v: a new simple path 1 ... u v
        count++;
    }
    cout << min(count, CAP) << "\n";
}
\end{lstlisting}
Samples: \textbf{3}, \textbf{16}, \textbf{2023}. Large tests: $K_{11}$ (9\,864\,101 paths) prints $1\,000\,000$ in 0.18\,s; the path on $2\cdot10^5$ vertices prints $200\,000$; a random graph with $N = M = 2\cdot10^5$ and all degrees $\le 10$ prints $1\,000\,000$ in 0.9\,s, of which 0.55\,s is reading the input. Checked against a bitmask DP ($\text{dp}[S][v]$ = number of simple paths from 1 with vertex set $S$ ending at $v$, summed and capped) on 500 random graphs with $N \le 11$: identical.
\end{solutionbox}

\begin{solutionbox}[E4.5 Unique Color]
The path from 1 to $x$ is exactly the DFS stack when $x$ is entered. Keep $\texttt{cnt}[c]$ = how often colour $c$ occurs on the stack: $x$ is good iff $\texttt{cnt}[C_x] = 0$ just before $x$ is pushed. Increment on enter, decrement on exit --- template 3 of the pattern. A tree has one path per vertex, so this is linear; the tree can be a path of $10^5$ vertices, hence the iterative DFS.
\begin{lstlisting}[style=cpp]
#include <bits/stdc++.h>
using namespace std;

// Unique Color (ABC198 E): x is good iff no vertex on the path 1..parent(x) has colour C[x].
// Path stack: cnt[c] = occurrences of colour c on the current root path. Iterative (n = 1e5).
int main() {
    ios::sync_with_stdio(false);
    cin.tie(nullptr);
    int n;
    cin >> n;
    vector<int> c(n + 1);
    for (int i = 1; i <= n; i++) cin >> c[i];
    vector<vector<int>> adj(n + 1);
    for (int i = 0; i < n - 1; i++) {
        int a, b;
        cin >> a >> b;
        adj[a].push_back(b);
        adj[b].push_back(a);
    }
    vector<int> cnt(100001, 0), par(n + 1, 0);
    vector<size_t> it(n + 1, 0);
    vector<char> good(n + 1, 0);
    vector<int> st(1, 1);
    good[1] = 1; cnt[c[1]]++;                       // enter the root
    while (!st.empty()) {
        int u = st.back();
        if (it[u] == adj[u].size()) {               // leave u: its colour leaves the path
            cnt[c[u]]--;
            st.pop_back();
            continue;
        }
        int v = adj[u][it[u]++];
        if (v == par[u]) continue;
        par[v] = u;
        good[v] = (cnt[c[v]] == 0);                 // the path 1..u is exactly the stack
        cnt[c[v]]++;                                // enter v
        st.push_back(v);
    }
    for (int i = 1; i <= n; i++)
        if (good[i]) cout << i << "\n";
}
\end{lstlisting}
Samples: \textbf{1 2 3 4 6} and \textbf{1 2 3 5 6 7 8}. A path of $10^5$ vertices runs in 0.5\,s (mostly input). Checked against a brute force (BFS parents, then walk up from every $x$ and compare colours, \bigO{n^2}) on 500 random trees with $N \le 12$ and 3 colours: identical.
\end{solutionbox}

\begin{solutionbox}[E4.6 Path Prefixes]
Keep on the stack the prefix sums of $b$ along the root path: $\texttt{pre}[0] = 0$, $\texttt{pre}[d]$ = sum of $b$ on the first $d$ edges. Entering $v$ pushes $\texttt{pre.back()} + b_v$, leaving $v$ pops it. Because all $b_j \ge 1$, \texttt{pre} is strictly increasing, so $r_v$ = (the largest $d$ with $\texttt{pre}[d] \le A_v$) is one \texttt{upper\_bound}. Total \bigO{n \log n}; the tree can be a chain of $2\cdot10^5$ vertices, so the DFS is iterative.
\begin{center}
\begin{tikzpicture}[x=1.05cm, y=0.8cm]
    \tikzset{gv/.append style={minimum size=6.2mm, font=\scriptsize\bfseries}}
    \node[gvseen] (p1) at (0,0) {1}; \node[gvseen] (p2) at (1.8,0) {2}; \node[gvseen] (p4) at (3.6,0) {4}; \node[gvcur] (p3) at (5.4,0) {3};
    \draw[gtree] (p1) -- node[gw, above]{5 / 6} (p2);
    \draw[gtree] (p2) -- node[gw, above]{9 / 10} (p4);
    \draw[gtree] (p4) -- node[gw, above]{5 / 1} (p3);
    \node[font=\scriptsize, anchor=west] at (6.0,0) {$A_3 = 5 + 9 + 5 = 19$};
    \foreach \v/\i in {0/0, 6/1, 16/2, 17/3} {
        \node[cell] (q\i) at (\i*1.8,-1.2) {\v};
        \node[idx] at (\i*1.8,-1.85) {\texttt{pre[\i]}};
    }
    \node[font=\scriptsize, anchor=west, align=left] at (6.0,-1.2) {largest $d$ with $\texttt{pre}[d] \le 19$:\\$d = 3$, so $r_3 = 3$};
\end{tikzpicture}
\end{center}
{\small Sample 1 at vertex 3 (edge labels $a / b$): the stack is the root path $1, 2, 4, 3$, and \texttt{pre} holds its $b$-prefix sums.}
\begin{lstlisting}[style=cpp]
#include <bits/stdc++.h>
using namespace std;
typedef long long ll;

// Path Prefixes (CF 1714G). Keep the prefix sums of b along the root path on a stack:
// pre[d] = b-sum of the first d edges of the path. Since b > 0, pre is increasing, so
// r_i = (largest d with pre[d] <= A_i) = upper_bound(pre, A_i) - 1. Iterative DFS.
int main() {
    ios::sync_with_stdio(false);
    cin.tie(nullptr);
    int t;
    cin >> t;
    while (t--) {
        int n;
        cin >> n;
        vector<vector<int>> child(n + 1);
        vector<ll> a(n + 1), b(n + 1);
        for (int j = 2; j <= n; j++) {
            int p;
            cin >> p >> a[j] >> b[j];
            child[p].push_back(j);
        }
        vector<ll> r(n + 1, 0), A(n + 1, 0);
        vector<ll> pre(1, 0);                    // the path stack: pre[0] = 0 for the empty prefix
        vector<size_t> it(n + 1, 0);
        vector<int> st(1, 1);
        while (!st.empty()) {
            int u = st.back();
            if (it[u] == child[u].size()) {      // leave u: its edge leaves the path
                st.pop_back();
                pre.pop_back();
                continue;
            }
            int v = child[u][it[u]++];
            A[v] = A[u] + a[v];
            pre.push_back(pre.back() + b[v]);    // enter v: push the b-sum of the path 1..v
            r[v] = (upper_bound(pre.begin(), pre.end(), A[v]) - pre.begin()) - 1;
            st.push_back(v);
        }
        for (int i = 2; i <= n; i++) cout << r[i] << " \n"[i == n];
    }
}
\end{lstlisting}
The sample (4 test cases) is reproduced: \texttt{0 3 1 2 1 1 2 3}, \texttt{0 0 3}, \texttt{1 2 2}, \texttt{0 1 2 1 1 2 2 1 1}. (When the root is popped, \texttt{pre} becomes empty --- harmless, the loop ends.) A chain of $2\cdot10^5$ vertices runs in 1.2\,s here, 0.74\,s of it reading the input. Checked against a brute force (walk every root path explicitly, \bigO{n^2}) on 500 random multi-test inputs with $n \le 10$: identical.
\end{solutionbox}

\begin{solutionbox}[E4.7 Grid Paths]
A plain backtracking that only respects the pattern is far too slow (the free squares have to be avoided wisely). Three cuts, all provable:
\begin{itemize}[leftmargin=*]
    \item \textbf{Target too early.} Standing on the lower-left square before step 48 is a failure: return.
    \item \textbf{Split.} If the square straight ahead (in the direction of the last move) is visited or outside, while the two side squares are free, return: the path so far together with the border is a closed curve through the current square that separates the two side squares, and both must still be visited.
    \item \textbf{Forced square.} A free neighbour $w$ (not the target) with at most one free neighbour of its own can only be entered from here, now --- later it would be a dead end. If two neighbours are forced, return; if one is, take it.
\end{itemize}
\begin{center}
\begin{tikzpicture}[x=0.55cm, y=0.55cm]
    \tikzset{c/.style={gcell, minimum size=5.3mm, font=\tiny}, v/.style={c, fill=hitfill}, cu/.style={c, fill=currentfill, draw=accentorange, line width=0.9pt},
             fr/.style={c, fill=basefill}}
    \begin{scope}
        \foreach \x in {0,...,6} \foreach \y in {0,...,6} \node[c] at (\x,-\y) {};
        \foreach \x in {0,...,3} \node[v] at (\x,0) {};
        \foreach \y in {1,...,5} \node[v] at (3,-\y) {};
        \node[cu] at (3,-6) {}; \node[fr] at (2,-6) {L}; \node[fr] at (4,-6) {R};
        \draw[warningred, line width=1.5pt] (2.5,-6.55) -- (3.5,-6.55);
        \node[font=\scriptsize\bfseries] at (3,1.0) {prune 1: moving down into the border};
        \node[font=\scriptsize, align=center, text width=4.4cm] at (3,-7.6) {L and R lie on different sides of the path: whichever way it turns, the other side is lost};
    \end{scope}
    \begin{scope}[xshift=5.6cm]
        \foreach \x in {0,...,6} \foreach \y in {0,...,6} \node[c] at (\x,-\y) {};
        \foreach \x in {0,...,6} \node[v] at (\x,0) {};
        \foreach \x in {2,...,6} \node[v] at (\x,-1) {};
        \node[cu] at (1,-1) {}; \node[fr] at (0,-1) {w};
        \node[font=\scriptsize\bfseries] at (3,1.0) {prune 2: a forced square};
        \node[font=\scriptsize, align=center, text width=4.4cm] at (3,-7.6) {$w$ has one free neighbour left (below it): enter it now, or it becomes a dead end};
    \end{scope}
\end{tikzpicture}
\end{center}
\begin{lstlisting}[style=cpp]
#include <bits/stdc++.h>
using namespace std;

// Grid Paths (CSES 1625): paths of 48 steps from the upper-left to the lower-left square of a
// 7x7 grid that visit every square once and match a pattern of D, U, L, R and '?'.
const int K = 7, STEPS = K * K - 1;
bool vis[K + 2][K + 2];          // a frame of visited squares around the grid: no bounds checks
string pat;
long long answer = 0;
const int dr[4] = {1, -1, 0, 0}, dc[4] = {0, 0, -1, 1};   // D U L R
const string DIRS = "DULR";

bool freeSq(int r, int c) { return !vis[r][c]; }

void go(int r, int c, int step, int lastDir) {
    if (r == K && c == 1) {                      // at the target (row K, column 1 in framed coordinates)
        if (step == STEPS) answer++;             // ... a full path; otherwise the target came too early
        return;
    }
    if (step == STEPS) return;
    // prune 1: the square straight ahead is blocked, but both side squares are free.
    // The free squares then fall into two parts, and the path can never visit both.
    if (lastDir >= 0) {
        int ar = r + dr[lastDir], ac = c + dc[lastDir];
        if (!freeSq(ar, ac)) {
            if (lastDir < 2) {                   // moving vertically: the sides are left and right
                if (freeSq(r, c - 1) && freeSq(r, c + 1)) return;
            } else {                             // moving horizontally: the sides are up and down
                if (freeSq(r - 1, c) && freeSq(r + 1, c)) return;
            }
        }
    }
    // prune 2: a free neighbour w (not the target) with at most one free neighbour of its own
    // must be entered now, from here; later it could only be a dead end. Two such squares: fail.
    int forced = -1;
    for (int d = 0; d < 4; d++) {
        int wr = r + dr[d], wc = c + dc[d];
        if (vis[wr][wc] || (wr == K && wc == 1)) continue;
        int freeNb = 0;
        for (int e = 0; e < 4; e++) freeNb += freeSq(wr + dr[e], wc + dc[e]);
        if (freeNb <= 1) {
            if (forced >= 0) return;
            forced = d;
        }
    }
    for (int d = 0; d < 4; d++) {
        if (forced >= 0 && d != forced) continue;
        if (pat[step] != '?' && pat[step] != DIRS[d]) continue;
        int nr = r + dr[d], nc = c + dc[d];
        if (vis[nr][nc]) continue;
        vis[nr][nc] = true;                      // enter
        go(nr, nc, step + 1, d);
        vis[nr][nc] = false;                     // leave
    }
}

int main() {
    cin >> pat;
    for (int i = 0; i < K + 2; i++) vis[0][i] = vis[K + 1][i] = vis[i][0] = vis[i][K + 1] = true;
    vis[1][1] = true;
    go(1, 1, 0, -1);
    cout << answer << "\n";
}
\end{lstlisting}
\begin{center}
\small
\renewcommand{\arraystretch}{1.2}
\begin{tabular}{@{}lrrr@{}}
\toprule
\textbf{input} & \textbf{answer} & \textbf{calls, prune 1 only} & \textbf{calls, prunes 1 + 2} \\
\midrule
the sample \texttt{??????R??????U???\dots LD????D?} & 201 & 2\,642\,437 & 117\,844 \\
48 times \texttt{?} & 88\,418 & 46\,816\,433 (1.0\,s) & 2\,866\,561 (0.22\,s) \\
\bottomrule
\end{tabular}
\end{center}
With prune 1 alone the all-\texttt{?} input takes a full second here --- exactly the time limit; prune 2 cuts the calls by a factor 16. Checked by running the same program with $K = 5$ and $K = 6$ (all-\texttt{?}: 86 and 1770 paths) against a version without any pruning, on 300 and 100 random patterns (each character \texttt{?} with probability 0.8): identical.
\end{solutionbox}
